\documentclass[10pt,a4paper]{amsart}
\usepackage[margin=19mm]{geometry}
\usepackage{amsmath,amssymb,mathtools}
\usepackage[scr=rsfs,cal=euler]{mathalfa}
\usepackage{xfrac}
\usepackage[unicode,hidelinks,bookmarks=false]{hyperref}
\newcommand{\ie}{i.e.\ }
\newcommand{\cf}{cf.\ }
\newcommand{\resp}{respectively\ }
\newcommand{\Subsection}{\subsection}
\providecommand{\nobreakdash}{\nobreak\hbox{-}}
\setlength{\emergencystretch}{2em}
\multlinegap0pt
\usepackage{algorithm}\usepackage{algpseudocode}
\usepackage{tikz}
\newtheorem{proposition}{Proposition}
\newcommand{\qb}[1]{\ensuremath{\!\left[{#1}\right]}}
\newcommand{\qp}[1]{\ensuremath{\!\left({#1}\right)}}
\title{Batched Wavefront Sweeps for Linear Transport Uncertainty Quantification}
\author{Tristan Pryer}
\date{Source: arXiv:2609.14496v1 (CC BY 4.0)}
\begin{document}
\maketitle

\noindent\emph{Source and licence.} Batched Wavefront Sweeps for Linear Transport Uncertainty Quantification. Version arXiv:2609.14496v1 (CC BY 4.0), distributed by arXiv under the Creative Commons Attribution 4.0 International licence, http://creativecommons.org/licenses/by/4.0/, as recorded in the arXiv licence metadata for this exact version and shown on the abstract page https://arxiv.org/abs/2609.14496v1. Full licence notice: \texttt{licenses/arxiv-2609.14496-CC-BY-4.0.txt}.\par\smallskip

\noindent\textit{Editorial scope.} Counted unit: the article's proposition prop:uq\_reverse\_mode\_wavefront\_adjoint (source0.tex lines 2055--2099). The article's complete original proof of it is delivered with it (source0.tex lines 2101--2124, one \texttt{proof} environment of 24 lines). No mathematics is written, repaired or re-derived: the statement and the proof are the article's own lines, in the article's own notation, and no retained line is substituted or re-flowed.  Retained uncounted as the source-local context the proof needs: lines 2041--2051.  Nothing was omitted from any retained span, so the package carries no omission record.  No retained line contains a citation, so the wrapper carries no bibliography and no bibliography entry.  No retained line contains a figure environment, an image-inclusion command or drawing code, so no image is delivered and the package declares no asset dependency.  Editorial formatting only, in addition: an AMS \texttt{amsart} preamble that restates the article's own declarations and macros used by the retained lines, namely the environments \texttt{proposition} on separate counters, together with the standard packages those lines need.  The retained environments print in this document as Proposition 1 in order of appearance, whereas the article prints the same items with the numbering of its own full document.  The complete native source of the article as distributed in the arXiv e-print for this exact version is bundled unchanged as \texttt{sources/210321/source0.tex}.
\begin{equation}
  \delta U_{\alpha,K}
  =
  A_{\alpha,K}^{-1}
  \qb{
    \delta b_{\alpha,K}
    -
    (\delta A_{\alpha,K})U_{\alpha,K}
  }.
  \label{eq:uq_local_solve_derivative}
\end{equation}

\begin{proposition}[Reverse-mode derivative on a fixed face-sign branch]
\label{prop:uq_reverse_mode_wavefront_adjoint}
Consider one fixed face-sign branch with fixed mesh topology, fixed
upwind graph and a fixed chosen topological cell ordering. Assume that
all local blocks, coupling blocks, source vectors and inflow vectors
are continuously differentiable functions of $\xi$ on this branch, and
that every local diagonal block is nonsingular there. Let
\begin{equation}
  L_h(\xi)U_h(\xi)=F_h(\xi)
  \label{eq:uq_global_triangular_system}
\end{equation}
denote the assembled upwind DG system ordered by the sweep, and let
$\mathcal{A}_h(\xi)$ denote the wavefront algorithm. Then
\begin{equation}
  \mathcal{A}_h(\xi)=U_h(\xi)=L_h(\xi)^{-1}F_h(\xi).
\end{equation}
Moreover, for any differentiable scalar discrete observable
$J_h(U_h(\xi),\xi)$, reverse-mode differentiation through the unrolled
wavefront algorithm gives
\begin{equation}
  D_\xi J_h(U_h(\xi),\xi)[\delta\xi]
  =
  \partial_\xi J_h(U_h(\xi),\xi)[\delta\xi]
  +
  \lambda_h(\xi)^T
  \qb{
    D_\xi F_h(\xi)[\delta\xi]
    -
    D_\xi L_h(\xi)[\delta\xi]U_h(\xi)
  },
  \label{eq:uq_wavefront_reverse_directional}
\end{equation}
where $\lambda_h(\xi)$ solves
\begin{equation}
  L_h(\xi)^T\lambda_h(\xi)
  =
  \qp{
    \frac{\partial J_h}{\partial U_h}(U_h(\xi),\xi)
  }^T.
  \label{eq:uq_wavefront_adjoint_system}
\end{equation}
For batched computations with several sweep classes, the identity holds
on each class and scalar observable derivatives are obtained by summing
the corresponding class contributions.
\end{proposition}

\begin{proof}
On a fixed branch, the wavefront algorithm is block forward
substitution for the ordered triangular system
\eqref{eq:uq_global_triangular_system}. Hence
$\mathcal{A}_h(\xi)=L_h(\xi)^{-1}F_h(\xi)$.

Differentiating \eqref{eq:uq_global_triangular_system} in direction
$\delta\xi$ gives
\begin{equation}
  L_h(\xi)\delta U_h
  =
  D_\xi F_h(\xi)[\delta\xi]
  -
  D_\xi L_h(\xi)[\delta\xi]U_h(\xi).
  \label{eq:uq_wavefront_forward_directional}
\end{equation}
Combining this with the observable derivative and the adjoint equation
\eqref{eq:uq_wavefront_adjoint_system} gives
\eqref{eq:uq_wavefront_reverse_directional}. Reverse-mode automatic
differentiation traverses the local operations in reverse topological
order. The transpose of each local solve is the transpose local block
solve associated with \eqref{eq:uq_local_solve_derivative}; the reverse
pass is therefore block backward substitution with $L_h(\xi)^T$.
\end{proof}

\end{document}
